%!TEX root = ../main.tex

\chapter{融合节点分类与图推理的金融反欺诈系统设计与实现}

在数字金融生态持续扩张与业务链路日益复杂的背景下，金融欺诈呈现出更强的对抗性、隐蔽性与团伙化特征，
其行为往往跨账户、跨设备、跨场景协同发生，并通过长链路交互关系掩盖真实风险来源。
在此情形下，依赖独立同分布假设的传统统计模型，或仅针对局部邻域进行判别的孤立节点分类方法，
难以稳定捕获跨实体、多跳关联的结构化攻击模式，也难以向风控决策提供可核验的证据链支撑。

为验证本文提出的结合强化学习推理微调的图结构基座模型在垂直风控领域的有效性与泛化能力，
本章设计并实现一套面向金融反欺诈的科研原型系统。
系统以统一的图表征为核心，将\emph{节点分类的风险感知}与\emph{图推理的风险归因}贯通起来，
形成从风险筛查、证据聚合到可解释结论输出的端到端闭环，
从而为复杂金融网络中的反欺诈研判与事后审计提供可解释、可追溯的技术支撑。

\section{金融反欺诈场景需求分析}
\textbf{1)\ 业务需求分析}

面向金融反欺诈场景，具体业务需求问题拆解如\autoref{fig:Financial_Requirements} 所示。
面向金融欺诈团伙化与策略快速演变的现实特征，系统应具备超越传统规则引擎的泛化识别能力，
利用图拓扑与跨实体关联线索，定位高隐蔽性风险主体，
并在样本稀缺或黑产作弊手法发生变化时保持模型与系统准确率和召回率的保持。
其次，业务关注的不仅是个体风险判别，更是风险在群体结构中的传播与组织化协同机制；
因此系统需能够在异构交互网络中挖掘隐性关联与团伙拓扑指纹，实现对欺诈链条的结构化刻画与批量化打击。
最后，考虑到风控合规审计与人工复核的刚性要求，系统必须提供可解释、可追溯的决策依据：
在给出风险评分与等级的同时，能够明确呈现关键关联邻居、风险传导路径与支撑证据，并以结构化证据链与自然语言解释的形式降低复核成本。


\textbf{2)\ 功能需求分析}

在设计融合节点分类与图推理的金融反欺诈系统时，其主要的功能需求如下：
首先，系统应当首先提供宏观的风险态势感知能力，允许快速识别高风险群体；
同时，必须支持对微观个体、交易链路及子图结构的像素级穿透分析，确保从“面”到“点”的逻辑连续性。
其次，针对图结构基座模型的黑盒特性，
可视化界面需将隐式的推理路径以及中间层的融合评分显性化，
将其转化为图谱上的高亮路径、关键节点标记以及自然语言解释报告。
另外，系统不应是单向的信息展示终端，而应是人机协作的接口。系统需要复核结论、标签修正以及对推理逻辑的反馈，
应能够实时回流至系统，用于后续的模型微调与策略迭代。

\begin{figure}[htbp]
    \centering
    \includegraphics[width=.8\textwidth]{5_system/Financial_Anti_Fraud_System_Biz_Requirements.pdf}
    \caption{融合节点分类与图推理的金融反欺诈系统业务需求图}
    \label{fig:Financial_Requirements}
\end{figure}

\textbf{3)\ 性能与可用性指标}

基于科研原型定位，性能指标以可用交互与可复现实验为约束，
侧重计算准确性、吞吐能力与结果一致性。在离线批处理层面，系统应在单机多卡环境下支持千级节点、万级边规模的中等图数据构建与评分，
并在分钟级时间内完成全量风险覆盖，以保证评测与演示的可操作性。在交互推理层面，面向专家审核的单次图推理端到端延迟应稳定在秒级（例如 $<5$s），
以维持分析流程的连续性。与此同时，系统需建立严格的版本化与复现机制，确保在相同输入与配置下风险评分与推理路径输出具有可追溯的一致性，
从而满足学术验证的严谨性要求。


\section{总体系统架构设计}
为实现上述业务目标，本系统自底向上形成从原始数据到可解释风险决策的端到端闭环。
该架构遵循数据流向的递进关系：下层负责将多源业务日志抽象为可计算的动态图结构，
中间层以节点分类完成全量风险筛查并产出高风险候选，上层以图推理生成可核验的证据与报告，
最终通过可视化与导出能力支撑人工复核与策略迭代。整体如\autoref{fig:Fraud_System_Functional_Architecture} 所示。

\begin{figure}[htbp]
    \centering
    \includegraphics[width=.7\textwidth]{5_system/Fraud_System_Functional_Architecture.pdf}
    \caption{融合节点分类与图推理的金融反欺诈系统功能架构图}
    \label{fig:Fraud_System_Functional_Architecture}
\end{figure}

\textbf{1）数据接入与图构建层：}
本层面向多源离线数据提供统一接入能力，覆盖 CSV、JSON 与交易流水等输入形式，并提供缺失处理、异常裁剪与时间窗聚合等预处理流程，以提升数据可用性与口径一致性。在此基础上，本层通过图结构映射组件将用户、设备、IP 等实体及其交互关系编码为异构图，并进一步构建按时间窗口切分的动态图快照，用于保留行为演化与跨实体关联结构。本层的目标是将复杂业务事件转化为可复现、可版本化的图数据输入，为后续模型推断提供稳定的数据底座。

\textbf{2）分类感知层：}
本层以通用图基座模型为核心，面向目标节点执行高通量风险评分，并完成风险等级划分，承担系统的前置过滤与候选筛选职责。该层输出两类结果：一是面向全量实体的可排序风险清单，支撑批处理筛查；二是供上游推理模块使用的高风险候选集合，用于将计算资源集中于最有价值的审查对象。本层的目标是在保证覆盖面的前提下实现快速定位与资源分配，为深度归因建立可控的输入范围。

\textbf{3）风险推理归因层：}
本层以强化学习推理模型为核心，围绕候选高风险实体开展结构化归因，生成具备逻辑一致性的推理链路，并同步完成关键风险路径提取与异常子图模式识别。在输出形态上，本层形成两类交付物：其一为可回溯的证据链，用于支撑审计与复核；其二为面向业务人员的自然语言风险报告，用于降低解读成本并提升处置效率。本层的目标是将判别式评分转化为可解释的关联证据，实现从风险发现到风险成因的闭环连接。

\textbf{4）全链路回溯与可视化层：}
本层面向分析与复核场景提供交互式图谱探索能力，支持多跳邻域展开与关键关系高亮，便于对风险链路进行逐步核验与定位。同时，本层提供风险清单导出与推理证据摘要沉淀能力，并支持离线复盘与策略迭代所需的结果留痕与对照管理。本层的目标是将模型输出转化为可操作的分析界面与可复现的实验记录，支撑人工审核闭环与论文实验复现。

% 你希望动态图快照的时间窗在论文中设定为多长，例如 1 小时、1 天或 7 天？

\begin{figure}[htbp]
    \centering
    \includegraphics[width=.7\textwidth]{5_system/Flowchart.pdf}
    \caption{融合节点分类与图推理的金融反欺诈系统整体流程图}
    \label{fig:Flowchart}
\end{figure}

结合\autoref{fig:Flowchart} 的分层功能栈与数据流向，
原型系统采用“判别式全量筛查—生成式深度归因”的两阶段闭环。
其核心思想是先用高吞吐判别模型完成全量风险分流，再仅对高风险少量样本调用生成式推理进行可解释归因。

具体流程如下：从原始日志出发，系统对交易与行为日志执行缺失修复、异常裁剪、时间窗聚合与统计特征抽取，得到结构化的实体与交互记录。
基于清洗后的实体与交互记录，系统按统一 Schema 将数据映射为异构图。
明确节点类型和边类型并固化字段、方向与约束，使后续模型在训练与推断时共享同一图输入规范。

在阶段一的全量风险感知中，
系统调用第~3~章图节点分类基座模型对目标节点进行批量风险评分，并以阈值规则完成风险分流：
当风险评分低于阈值时，样本直接判定为低风险并通过；当风险评分超过阈值时，样本被标记为高风险候选并进入阶段二。
该阶段的输出是可排序的高风险候选集合，以最小化阶段二的计算负载并保证全量覆盖。

在阶段二的深度风险归因中，系统围绕高风险候选先执行目标子图提取，
构造与当前决策相关的上下文子图作为推理输入；随后调用第~4~章强化学习图推理模型生成逻辑一致的推理链路，
并输出风险定性结论与结构化证据链。

最终，系统将证据链与风险要点映射到可视化研判界面，支持多跳邻域展开、关键关系高亮与报告生成，
实现模型输出的可核验呈现与可复盘沉淀，形成闭环工作流。


\section{金融风控图数据构建与推理样本生成}
本节旨在详细阐述如何将原始的多源金融业务数据，转化为图基座模型可理解、可推理的高质量序列化输入。
这一过程涵盖了从底层的图谱本体设计、多源异构数据的清洗与动态图构建，到适配欧拉路径编码的序列化处理，
以及基于强化学习需求的高质量思维链推理样本生成的全生命周期。
为第5.4节的模型训练与第5.6节的系统测评奠定坚实的数据基础。

\textbf{1)金融异构图谱本体设计与多模态属性建模}

不同于通用知识图谱，风控图谱的设计需紧密围绕风险定义展开，强调实体属性的区分度与关系类型的因果性。
基于对真实信贷与支付风控业务的深入拆解，本文构建了一个包含5类核心节点与6类关键边关系的异构信息网络。
具体的数据项定义与处理方式如\autoref{tab:feature_processing_cn}所示。


\textbf{节点类型定义}： 在金融网络中，节点是风险的载体。本文不仅关注账户本身，更关注承载账户的物理媒介与环境。

\begin{enumerate}
    \item 用户： 交易的主体。
    包含注册ID、年龄段、职业的静态身份画像与动态风险标签。

    \item 设备：设备指纹唯一标识，承载了操作系统、
    模拟器特征、越狱状态等关键环境属性。

    \item IP地址：网络环境标识。包含地理归属、运营商类型等属性，
    用于地理位置异常。

    \item 商户：资金流转的汇聚点。包含商户类别码、经营主体信息，用于识别虚假交易。

    \item 关联实体：如MAC地址、加密的身份证号、手机号等，构建跨账户的隐性关联。
\end{enumerate}

\textbf{边关系类型定义：}边承载了风险传播的动态语义。

\begin{enumerate}
    \item 资金交易： 有向边。携带金额、时间、附言等核心属性。

    \item 设备登录： 无向边。高频切换设备往往意味着盗号或代办风险。

    \item 网络接入： 无向边。用于发现聚集性风险，如同一网段下的批量注册。

    \item 社交关联： 无向边。基于转账频繁联系人、通讯录或担保关系构建。

    \item 物理共现： 无向边。用于挖掘线下聚集的黑产窝点。
    \item 属性共享： 无向边。如共用手机号、共用身份证等团伙强特征。
\end{enumerate}

\begin{table*}[h]
\centering
\small
\caption{金融反欺诈系统中图谱多模态属性定义与处理规范}
\label{tab:feature_processing_cn}
\renewcommand{\arraystretch}{1} % 适当增加行高，中文排版更美观
\begin{tabular}{p{0.1\linewidth} p{0.1\linewidth} p{0.2\linewidth} p{0.4\linewidth}}
\toprule
\textbf{节点/边} & \textbf{属性域} & \textbf{代表性数据项} & \textbf{构建与处理策略} \\
\midrule

% 用户节点部分
\multirow{6}{*}{\textbf{用户节点}} 
& 身份画像 & \texttt{user\_id}, \texttt{age\_group}, \texttt{occupation} & \textbf{ID映射与离散化}：结构化Token映射为唯一ID；时间特征离散化处理；职业信息映射为文本描述。 \\
\cmidrule(l){2-4}
& 信用历史 & \texttt{credit\_score}, \texttt{overdue\_count} & \textbf{分位数分箱}：将连续的信用评分、逾期次数等数值通过分位数划分为不同的Token区间。 \\
\cmidrule(l){2-4}
& 文本备注 & \texttt{reg\_memo}, \texttt{complaint\_history} & \textbf{语义嵌入}：保留原始自然语言文本，经由RoBERTa编码器提取语义向量，避免One-hot造成的稀疏性。 \\
\midrule

% 设备节点部分
\multirow{3}{*}{\textbf{设备节点}} 
& 环境指纹 & \texttt{device\_id}, \texttt{is\_root}, \texttt{battery\_status} & \textbf{结构化映射}：布尔值直接映射；电量等动态特征进行离散化处理。 \\
\cmidrule(l){2-4}
& 异常特征 & \texttt{ua\_string}, \texttt{hook\_log} & \textbf{原始序列保留}：保留User-Agent字符串原文，用于大语言模型(LLM)识别特定模拟器特征。 \\
\midrule

% 交易边部分
\multirow{3}{*}{\textbf{交易边}} 
& 核心流向 & \texttt{amount}, \texttt{timestamp}, \texttt{currency} & \textbf{数值变换与时序嵌入}：金额进行对数缩放(Log-scaling)；计算时间戳差值 $\Delta t$ 进行嵌入。 \\
\cmidrule(l){2-4}
& 语义描述 & \texttt{tx\_memo}, \texttt{product\_name} & \textbf{文本拼接}：将交易附言、商品名称等直接作为边的文本属性进行拼接与编码。 \\
\midrule

% 登录边部分
\textbf{登录边} 
& 上下文环境 & \texttt{geo\_hash}, \texttt{auth\_method} & \textbf{混合编码}：将地理位置GeoHash编码为Token，用于后续的空间距离计算。 \\

\bottomrule
\end{tabular}
\end{table*}


\textbf{2)面向金融交易网络的欧拉路径图编码实现}

第3章提出的欧拉路径编码是本系统的核心技术支撑。
然而，通用图算法针对的是无向或简单有向图，
而金融交易网络是典型的强有向、多重边且属性密集的图结构。
本节详细阐述如何将该技术适配于金融场景，实现对资金流向的无损序列化。
具体实现而言，包含以下步骤：

\textbf{资金流向保真的欧拉化预处理}:欧拉路径存在的充要条件是
图中每个节点的入度等于出度或存在两个奇度数节点。
真实的交易子图往往是不平衡的。直接加虚边可能会混淆真实资金流与虚拟平衡流。
因此，本文引入以下适配改进：
首先引入全局虚拟节点$``$，单向连接至子图中的所有孤立分量核心节点，
确保图连通。
其次，针对度数不平衡节点，添加标记为 $``$ 的回溯虚边。
不同于通用算法的随机加边，本文优先在属于同一资金链路的节点间添加回溯边，
在序列化时， 边会被分配特殊的Mask Token，
使模型在自注意力计算时能够区分真实交易与补图路径，
既满足了欧拉路径的数学存在性，又保留了真实资金流向的单向语义。

\textbf{时序优先的深度搜索序列化:}在生成欧拉路径时，
节点遍历的顺序直接决定了序列的语义连贯性。针对金融场景，本文改进了路径搜索策略，
采用{时序优先的深度搜索}。具体而言：在节点 $u$ 选择下一条边 $(u, v)$ 进行遍历时，不进行随机选择，
    而是对所有出边按交易时间戳 $t$ 进行升序排列，优先遍历发生时间最早的交易。
    从而生成的结果序列 $ User_A -> Trans_1(t1) -> User_B -> Trans_2(t2) -> User_C $
    自然地蕴含了 $t1 < t2$ 的时序约束。这使得Transformer的位置编码
    能够隐式地捕捉到资金流动的因果链条，极大增强了模型对“快进快出”、“分层转账”模式的识别能力。

\textbf{全属性模式序列构建:}为了承载5.3.2节定义的多模态丰富属性，
系统采用第3章定义的全属性模式进行最终的Token序列构建。
每个边步长被展开为包含节点、边及其属性的复合序列。
具体序列结构示例如\autoref{tab:sequence_token_input}所示。
这种极其稠密的编码方式虽然增加了序列长度，但最大限度地保留了风控所需的全部细节，
使基座模型能够同时利用拓扑结构和语义特征进行联合推理。

\begin{table}[h]
\centering
\small
\caption{金融反欺诈系统中欧拉图遍历路径序列化Token输入示例}
\label{tab:sequence_token_input}
\renewcommand{\arraystretch}{1}
\begin{tabular}{c p{0.15\linewidth} p{0.15\linewidth} p{0.35\linewidth}}
\toprule
\textbf{位置索引} & \textbf{Token示例} & \textbf{语义角色} & \textbf{详细说明} \\
\midrule
$t$ & \texttt{UID\_258} & 当前节点ID & 源用户的唯一标识符，经由ID Mapping映射。 \\
\midrule
$t+1$ & \texttt{[NODE\_ATTR]} & 节点离散属性 & 聚合后的结构化特征，经离散化处理后的Token。 \\
\midrule
$t+2$ & \texttt{[SEM\_EMB]} & 文本语义向量 & 
用户备注信息({reg\_memo})经RoBERTa编码后的稠密向量嵌入。 \\
\midrule
$t+3$ & \texttt{[EDGE\_OUT]} & 边类型 & 显式指示关系方向为资金流出，区分出入度语义。 \\
\midrule
$t+4$ & \texttt{[EDGE\_ATTR]} & 边属性特征 & 离散化的交易特征组合
（如：{大额}、{短时}、{夜间}）。 \\
\midrule
$t+5$ & ``Payment for...'' & 交易附言 & 原始交易文本序列（保留自然语言语义，作为序列的一部分）。 \\
\midrule
$t+6$ & \texttt{UID\_996} & 下一跳节点ID & 目标交互对象（Counterparty），用于引导模型预测链路关系。 \\
\bottomrule
\end{tabular}
\end{table}




\textbf{3)风险推理样本的定义与生成体系}

为了支撑第4章所述的结合强化学习微调的图推理，
必须构建高质量、包含显式思维链的样本集。
在金融风控领域，这意味着模型不仅要输出“通过/拒绝”，
还要输出“因为A在短时间内从B处收到大额资金并立即转给高风险账户C，
且设备环境与历史习惯不符，判定为洗钱风险”的完整逻辑链条。

因此，本文将金融图推理任务形式化为三元组 $(G, q, y)$：
$G$表示为经欧拉路径序列化后的风控子图Token序列，
$q$是自然语言形式的风险排查指令以及包含思维链与答案的序列 $y = (c, a)$。
$c$是分步骤的推理过程，$a$代表结构化的结论，通常是JSON格式，便于下游系统解析。
针对图推理任务，本节的系统提示词覆盖三个方面，具体如下：
\begin{prompt}{风险感知类的风险排查指令}
请分析中心用户$User_X$是否存在异常，并给出风险评分（0-100）。
\end{prompt}
\begin{prompt}{归因解释类的风险排查指令}
$User_X$被标记为高风险，请列出3个关键交易和1个关联团伙，并说明理由。
\end{prompt}
\begin{prompt}{模式挖掘类的风险排查指令}
图中是否存在资金闭环或设备共享行为？请输出具体的节点路径$ID$。
\end{prompt}

\textbf{4)基于教师模型蒸馏的冷启动样本生成}

由于缺乏大规模的专家标注CoT数据，本系统采用
结合教师模型蒸馏和规则校验的策略生成冷启动样本库（参考章节 \autoref{subsec:grpo_train_detail}）。
教师模型配置选用推理能力较强的通用大模型Qwen2.5-32B-Instruct作为教师。
提示词工程设计包含领域知识的系统提示词：

\begin{prompt}{基于教师模型蒸馏的冷启动样本生成系统提示词}
你是一名资深金融反欺诈/反洗钱专家。
给定输入的“欧拉序列文本”（将图事件按欧拉路径/游走序列化），
你需要{还原图结构}并{判断是否命中典型欺诈拓扑}，生成可用于教师模型蒸馏的推理样本。
\end{prompt}



最后将构建好的子图序列 $G$ 和指令 $q$ 输入教师模型，
要求其按照对应步骤生成思维链结果。
为了确保生成的推理样本格式正确且符合事实，
本文利用第4章提出的可验证奖励函数 $\mathcal{R}(G, q, y)$ 
对生成的思维链数据进行严格过滤，选取高样本纳入数据集，并
利用拓扑一致性校验解析所有交易链路，
检查其是否真实存在于输入图 $G$ 中。若模型凭空捏造事实，则该样本直接剔除。
最后使用图算法辅助验证，检查推理过程中的定性与最终评分是否一致。

通过上述流水线，本文构建了包含 5,000条 高质量图推理指令样本的数据集，称为Fin-CoT-5K。
样本包含40\%风险感知问题、35\%归因解释问题以及25\%的模式挖掘问题，
覆盖了洗钱、赌博、电信诈骗、信贷中介、设备农场等5大类典型黑产场景。
经过人工抽检，生成样本的准确率达到98.5\%，逻辑连贯性达到95\%以上，
该数据集用于后续模型在反欺诈垂直场景的微调训练，并集成于系统中。

\section{金融风控大模型训练与优化}
%（如需子图排版，请确保导言区已包含：\usepackage{subcaption}）

\section{风控决策可视化与解释性模块}
\label{sec:vis_explain_module}

本节将重点阐述金融反欺诈系统的具体功能实现和可视化界面。
该模块作为连接底层算法引擎与上层业务功能交互的桥梁，
集成展示了从全量风险感知到深度推理归因的全过程，通过多视图协同实现对模型决策逻辑的白盒化呈现。
该模块是本原型系统的核心交互界面，也是本文所提出的图基座模型技术落地能力的最直接体现。

整体功能界面如\autoref{fig:initial_ui}所示。
整体功能架构划分为四个紧密耦合的子区域：

\begin{figure}[htbp]
    \centering  % 让图片在figure环境中居中显示
    \includegraphics[width=1\linewidth]{5_system/prototype/initial_ui.png}  % 图片宽度设为0.6倍行宽
    \caption{融合节点分类与图推理的金融反欺诈系统的整体初始界面}
    \label{fig:initial_ui}
\end{figure}

\begin{enumerate}
    \item {配置与控制区：} 
    负责数据接入、参数调优与分析任务调度。
    \item {全量风险感知区：} 
    展示基于节点分类模型的批量打分结果与风险分层。
    \item {图谱拓扑探索区：} 
    提供基于力导向布局的动态子图可视化，展示交易网络结构。
    \item 推理归因与报告区：
    展示基于强化学习微调模型生成的自然语言证据链、风险定性结论及处置建议。
\end{enumerate}

下文将结合系统实现的具体界面，详细阐述各子模块的设计细节与业务价值。


\textbf{1)全量数据接入与参数化配置}

如\autoref{fig:full_data_input_config}所示，系统支持标准化的CSV数据导入，
明确区分源节点、目标节点及其对应的属性列。
为了适配第3章所述的欧拉路径图编码器，
前端配置界面允许用户指定K-hop采样阶数与随机种子。
K-hop采样阶数控制子图构建的深度。
较小的 $K$ 值聚焦于直接交易对手，适用于快速筛查；
较大的 $K$ 值捕获长链路的洗钱路径，适用于深度挖掘。
系统设计了Mock生成子图与上传数据构建两种模式。
Mock模式利用内置的生成器实时构造典型的欺诈拓扑，
用于算法演示与基准测试；真实模式则直接对接业务数据流，支撑实战分析。

\begin{figure}[htbp]
    \centering
    \begin{subfigure}[t]{0.49\linewidth}
        \centering
        \includegraphics[width=\linewidth]{5_system/prototype/initial_input.png}
        \caption{数据上传与参数化配置界面}
    \end{subfigure}
    \hfill
    \begin{subfigure}[t]{0.49\linewidth}
        \centering
        \includegraphics[width=\linewidth]{5_system/prototype/initial_mork.png}
        \caption{Mock生成子图界面}
    \end{subfigure}
    \caption{金融反欺诈系统的全量数据接入与参数化配置界面}
    \label{fig:full_data_input_config}
\end{figure}

为适配业务实际需求，
为了解决风控中误杀与漏网的权衡问题，
系统在分析面板见\autoref{fig:threshold_setting}提供了可视化的阈值滑块组件。

\begin{figure}[htbp]
    \centering
    \includegraphics[width=0.6\linewidth]{5_system/prototype/threshold_setting.png}
    \caption{金融反欺诈系统的可视化阈值滑块组件}
    \label{fig:threshold_setting}
\end{figure}


\textbf{2)批量风险感知与分层筛选}

在完成数据接入后，系统首先调用第3章训练的节点分类图结构基座模型，
对输入账户列表进行全量风险评分。
这一过程对应系统架构分类感知层中的高通量风险评分和风险等级划分阶段。
如\autoref{fig:batch_score_output}所示，系统以结构化表格形式展示批量分析结果。表格包含以下关键字段：
\begin{itemize}
    \item Account：待检测账户ID，例如 $U_{1234}$;
    \item Prob：输出的初始欺诈概率（0-1区间），反映模型对该节点属于欺诈类别的置信度。
    \item Score：映射后的信用评分（0-100分制）。通常分数越高代表风险越高，
    \item Level：依据前述阈值配置自动划分的风险等级。
\end{itemize}

\begin{figure}[htbp]
    \centering
    \includegraphics[width=0.6\linewidth]
    {5_system/prototype/batch_score_output.png}
    \caption{金融反欺诈系统输出待检测账户的风险等级}
    \label{fig:batch_score_output}
\end{figure}

对于被标记为高风险的账户，
单一的评分数字不足以支撑封号冻结等强处置操作。
系统提供交互式图谱拓扑探索进行动态的交易结构可视化。
本模块的图谱展示和多跳探索区域如\autoref{fig:subgraph}所示。

\begin{figure}[htbp]
    \centering
    \includegraphics[width=0.6\linewidth]
    {5_system/prototype/subgraph.png}
    \caption{金融反欺诈系统的交互式图谱拓扑探索}
    \label{fig:subgraph}
\end{figure}

系统采用力导向算法对目标账户的K-hop子图进行渲染。
中心节点为蓝色：位于几何图的中心，表示待分析的目标账户。
异常节点为红色，表示为模型识别出的高危关联节点。
正常节点为灰色，用于还原完整的交易上下文。
设备媒介节点为紫色或者其他色，代表物理设备或IP环境。

系统同时实现了交易链路的可视化，
节点之间的连线代表交易或关联关系，
连线上的数值直接展示了交易金额或频次权重。箭头指示了资金流向。
在\autoref{fig:subgraph}中可以清晰地看到以 $U_{1234}$ 为核心的星状或发散状结构，
大量资金从中心节点流向四周（如 $U_{1234}$ -> $U_{4517}$, $U_{1234}$ -> $U_{2542}$）
。同时，右下角的红色设备节点$D_{5678}$通过虚线与中心节点相连，代表弱关联或设备共用关系。
这种视觉结构强有力地暗示了团伙作案或设备关联的风险模式。

为了防止可视化带来的信息丢失，
系统在图谱下方提供了原始数据浏览面板，具体如\autoref{fig:original_data_review}所示。
节点表展示每个节点的入度、出度、
是否中心节点及是否异常。
例如，$U_{9313}$ 的 is\_anomaly=1，与其在图谱中的红色显示对应。
边表展示每笔交易的源、目、类型、设备绑定、
金额及时间戳。 
这种图表联动的设计，
确保既能获得拓扑直觉，又能核实底层数据的准确性，符合可信AI的设计要求。

\begin{figure}[htbp]
    \centering
    \begin{subfigure}[t]{0.49\linewidth}
        \centering
        \includegraphics[width=\linewidth]{5_system/prototype/data_review_of_node.png}
        \caption{节点表原始数据浏览面板}
    \end{subfigure}
    \hfill
    \begin{subfigure}[t]{0.49\linewidth}
        \centering
        \includegraphics[width=\linewidth]{5_system/prototype/data_review_of_edge.png}
        \caption{边列表原始数据浏览面板}
    \end{subfigure}
    \caption{金融反欺诈系统的原始数据的透明化查阅}
    \label{fig:original_data_review}
\end{figure}

\textbf{3)基于强化学习微调的图基座模型的解释性推理与证据看板}

这是本系统最核心的创新模块，集成了第4章所述的强化学习微调图大模型。当点击某个高风险账户时，
系统会实时生成一份包含自然语言解释、风险评分及关键证据链的推理报告。

如\autoref{fig:original_data_review}所示，面板顶部展示了综合风险评分仪表盘，这里显示了风险评分和融合决策。
风险评分与上述介绍一致，KPI摘要则
列出了子图的关键统计指标，
如节点数、边数、异常点数、环路数等，
帮助使用者快速建立对案件复杂度的认知。
此外，系统利用微调后的大模型，生成了结构化的推理报告，对应
\autoref{fig:original_data_review}左侧文本区，包含三部分要素：（1）
结论先行：直接给出账户反欺诈风险等级的定性结论。
（2）可解释证据（Evidence）：报告列出了支撑结论的关键发现，
每一条发现都对应图谱中的具体子结构。
（3）关键路径征据：系统详细列出了具体的资金链路，
路径级的解释能够直接追踪资金去向。

\begin{figure}[htbp]
    \centering
    \begin{subfigure}[t]{0.49\linewidth}
        \centering
        \includegraphics[width=\linewidth]{5_system/prototype/evidence_low_score_case.png}
        \caption{低风险样例的风险评分与解释性推理}
    \end{subfigure}
    \hfill
    \begin{subfigure}[t]{0.49\linewidth}
        \centering
        \includegraphics[width=\linewidth]{5_system/prototype/evidence_mid_score_case.png}
        \caption{中等风险样例的风险评分与解释性推理}
    \end{subfigure}
    \caption{金融反欺诈系统中基于图推理模型的决策输出与证据摘要}
    \label{fig:original_data_review}
\end{figure}

最后，为了辅助对单一案件的判断，系统还提供了统计维度的辅助分析视图。
如\autoref{fig:amount_hist}所示，系统自动计算子图内所有交易金额的分布情况。通过
长尾分布识别直方图可以帮助识别是否存在大量小额测试单或单笔特大金额转账。
图中显示在2500-3000金额区间存在集中交易（频次为4），
这可能意味着欺诈团伙采用了固定额度的批量转账策略。

\begin{figure}[htbp]
    \centering
    \includegraphics[width=0.6\linewidth]
    {5_system/prototype/amount_hist.png}
    \caption{金融反欺诈系统的交易金额分布分析视图}
    \label{fig:amount_hist}
\end{figure}



\section{系统测评与场景验证}
\section{本章小结}

